%%%%%%%%%%%%%%%%%%%%%%%%%%%%%%%%%%%%%%%%%%%%%%%%%%%%%%%%%%%%%%%%%%%%%%%%%%%%%%
%%  Information sheet required with every JAAMAS submission.
%%  JAAMAS returns papers without review when this sheet is incomplete or
%%  uninformative, so each of the four questions is answered directly.
%%  Build:  pdflatex information_sheet
%%%%%%%%%%%%%%%%%%%%%%%%%%%%%%%%%%%%%%%%%%%%%%%%%%%%%%%%%%%%%%%%%%%%%%%%%%%%%%
\documentclass[10pt,a4paper]{article}
\usepackage[T1]{fontenc}
\usepackage{lmodern}          % scalable faces, needed by microtype expansion
\usepackage[margin=1.9cm]{geometry}
\usepackage{amsmath,amssymb}
\usepackage{parskip}
\usepackage[colorlinks=true,allcolors=blue!55!black]{hyperref}
\usepackage{xcolor}
\usepackage{microtype}
\pagestyle{plain}

\newcommand{\AS}{\textsf{AS}}
\newcommand{\AU}{\textsf{AU}}
\newcommand{\CS}{\textsf{CS}}
\newcommand{\CAS}{\textsf{CAS}}

\begin{document}
\small                       % keeps the sheet inside the two pages JAAMAS asks for

\begin{center}
{\large\bfseries Information sheet}\\[2pt]
{\bfseries Safe designs, unsafe ecosystems: deployment-layer selection in an
evolutionary AI race}\\[2pt]
Regular research paper
\end{center}

\section*{1. What is the main claim, and why does it matter to the autonomous
agents and multi-agent systems literature?}

In an ecosystem of deployed agents, the payoff an agent earns and the quantity
that decides whether its design is replicated are different functionals. The
agent plays; a principal who is not a player decides whether to keep deploying
it, and internalises only a fraction $\lambda$ of the harm the design causes
third parties. We call that gap the \emph{selection wedge}. Writing $a$ for the
expected task payoff of an interaction, $m$ for its expected number of unsafe
actions and $h$ for the harm each one causes, the population of designs evolves
under the replicator equation of the principal's functional
$\pi_{P}=a-\lambda h\,m$ rather than of the social one. That functional depends
on pass-through and harm only through the single product $L=\lambda h$, and it
makes every invasion condition affine in $L$. The claim is that governance of
such an ecosystem binds at the layer that selects designs and not at the layer
that plays them: an instrument applied to the agents can be exactly inert,
while the instrument that does move the vector field is not monotone in its own
strength.

This matters to the autonomous agents and multi-agent systems literature
because that literature builds, evaluates and deploys populations of agents
that principals choose, and every instrument it has for open agent populations
-- social laws, electronic institutions, norm synthesis and sanctioning, trust
and reputation, reward transfers and formal contracts -- acts from inside the
interaction and so moves the behaviour of a population whose membership is
given. Its evolutionary models of AI races, institutional incentives and
delegation identify earned payoff with reproductive fitness, using one payoff
matrix twice. Deployed ecosystems break that identification structurally, and
the two-matrix version yields guidance the one-matrix version cannot express:
which admission tests are dynamically empty and what to replace them with, why
a population that looks stationary can be losing its protection, and why an
accountability instrument has a range of strengths worse than a weaker
setting.

\section*{2. What is the evidence, and what does it imply?}

The interaction layer is a four-design race -- always safe $\AS$, always unsafe
$\AU$, conditionally safe $\CS$, conditionally antisocial $\CAS$ -- taken
unchanged from the framed AI-race experiment of Fern\'{a}ndez Domingos and Han
(preprint, arXiv:2607.26034, July 2026) and evaluated by exact expectation over
the horizon law rather than by sampling. An invader beats a resident when its
payoff and unsafe-action advantages satisfy $\delta a-L\,\delta m>0$, so each
of the twelve ordered invasions has a critical liability
$L^{*}=\delta a/\delta m$ wherever the two designs differ in unsafe count.
Four of these, together with one interior threshold, organise the paper and
span a factor of $368$: $0.116<0.551<4.274<42.765$. Each is a property of
$\pi_{P}$ alone and is therefore shared by every imitative payoff-monotone
dynamics, which we confirm for the logit, best-response and aspiration rules.

\emph{Solo evaluation is dynamically empty.} In this pool $\CAS$ weakly
dominates $\AU$ at every $L\ge0$, so a filter that scores a design against a
safe environment removes only a weakly dominated design; the attractor it
deletes carries the maximal long-run unsafe frequency $U=1$, and below
$L=0.116$ the surviving pool reproduces that same all-unsafe state on a
different design. Filtered and unfiltered long-run unsafe frequencies agree to
within $6.2\times10^{-3}$ on a common start measure. The dominance is a
property of the pool, not a theorem about filters: it fails once the pool
contains a design that punishes for longer than it is punished, which is the
one design-layer intervention the selection layer does not undo. Over all
fifteen probe sets the separating margin takes exactly two values, $0.539$ when
the probe reciprocates and contains no always-unsafe design and $0.132$
otherwise, so a reciprocating probe needs $7$ episodes at zero execution error
and $16$ at a $5\%$ rate where a solo probe needs $111$ and $136$. The
implication is that an evaluation protocol should be specified by who the agent
is evaluated against, not by how strict the threshold is.

\emph{Protection is governed by a variable that drifts.} $\AS$ and $\CS$ are
exact payoff twins, so the safe face carries no selection gradient, yet the
liability needed to repel an invader falls from $42.765$ at the unconditional
vertex to $0.551$ at the conditional one, a factor of $77.7$ at the baseline
calibration and between $4.4$ and $90.4$ across the robustness grid. Nothing
observable in the interaction changes as the population drifts along that face,
so resilience has to be monitored as a composition rather than as an outcome.

\emph{Long-run harm is not monotone in liability.} A mixture-free invasion
lemma, valid for any symmetric matrix game, gives the guard threshold
$L^{\dagger}=4.274$ exactly. On $(4.274,\,42.765)$ liability taxes the
retaliation that conditional safety depends on, the unsafe rest point becomes
uninvadable by conditional safety, and the basin-averaged unsafe frequency
jumps discontinuously from $0$ to $0.32$ as $L$ crosses the lower edge from
below. In a finite population the resulting trap is long-lived, with mean
escape times from $65$ generations at $Z=30,\beta=0.02$ to
$3.9\times10^{13}$ at $Z=100,\beta=0.2$. A dangerous range of this kind exists
at all $24$ prize and risk grid points, under all five charge schedules we
tried, and in every enlarged design pool at zero execution noise; its location
is not transferable, since the lower edge moves from $4.274$ in the reduced
pool to between $12$ and $13$ in the full one, and by an execution error rate
of $5\%$ the range has closed because the state liability was protecting has
itself stopped being safe. Assortative matching shrinks the range from above
and closes it at $r=0.354$, whereas misattributing a fraction $q$ of the blame
pushes the upper edge up exactly as $1/(1-q)$ while moving the lower edge by at
most $11.8\%$, so at $q=0.5$ the dangerous range has doubled. The implication
for a regulator is that raising liability out of a protected regime has to
clear the whole range in one step, and that incident records which book an
interaction against the wrong party widen the range they have to clear.

\emph{Losing a protected regime costs more than keeping it.} Coupling
enforcement to a monotone environmental state that erodes it to a residual
fraction $\rho$ makes the transition hysteretic: the ratio of the liability at
which the regime is recovered to the one at which it was lost is exactly
$1/\rho$, for every erosion channel and every time-scale separation, provided
the protected state is exactly safe. The cheap intervention is the early one.

\section*{3. Closest related contributions by other authors}

\textbf{Evolutionary models of AI races} are the closest work: Han, Pereira,
Santos and Lenaerts, \emph{Journal of Artificial Intelligence Research}
\textbf{69} (2020) 881--921; Cimpeanu et al., \emph{Scientific Reports}
\textbf{12} (2022) 1723; Alalawi et al., \emph{Applied Mathematics and
Computation} \textbf{508} (2026) 129627; surveyed in Han, \emph{AI
Communications} \textbf{35} (2022) 327--337. There the racing actor and the
reproducing actor coincide. We keep their race and replace the selection layer,
so the contribution is orthogonal rather than corrective.

\textbf{Evolutionary methods and normative design in multi-agent systems.}
Evolutionary dynamics are a standard tool of this journal's literature:
Bloembergen, Tuyls, Hennes and Kaisers, \emph{Journal of Artificial
Intelligence Research} \textbf{53} (2015) 659--697; Tuyls et al.,
\emph{Autonomous Agents and Multi-Agent Systems} \textbf{34} (2020) 7; Phelps,
McBurney and Parsons, \emph{Autonomous Agents and Multi-Agent Systems}
\textbf{21} (2010) 237--264. Norm and institution design supplies the
instruments we compare against: Shoham and Tennenholtz, \emph{Artificial
Intelligence} \textbf{73} (1995) 231--252; Morales, Wooldridge,
Rodr\'{i}guez-Aguilar and L\'{o}pez-S\'{a}nchez, \emph{Autonomous Agents and
Multi-Agent Systems} \textbf{32} (2018) 635--671; Morris-Martin, De Vos and
Padget, \emph{Autonomous Agents and Multi-Agent Systems} \textbf{33} (2019)
706--749. All of these act on agents inside the interaction; our instrument
acts on a principal outside it and changes which designs exist.

\textbf{Evaluation and reward modification in multi-agent settings.} Our
evaluation result speaks to Omidshafiei et al., \emph{Scientific Reports}
\textbf{9} (2019) 9937, which ranks agents by an evolutionary process over the
interaction rather than by a solo score, and our non-monotone instrument
contrasts with Willis, Du, Leibo and Luck, \emph{Autonomous Agents and
Multi-Agent Systems} \textbf{38} (2024) 49, and Haupt, Christoffersen, Damani
and Hadfield-Menell, \emph{Autonomous Agents and Multi-Agent Systems}
\textbf{38} (2024) 51, whose transfers and contracts act among the agents
themselves and admit a computable minimum. The closest architecture is Zheng,
Trott, Srinivasa, Parkes and Socher, \emph{Science Advances} \textbf{8} (2022)
eabk2607, whose outer planner optimises an instrument over a responding
population; ours only selects.

\textbf{The indirect evolutionary approach}: Alger and Weibull,
\emph{Econometrica} \textbf{81} (2013) 2269--2302, surveyed in \emph{Annual
Review of Economics} \textbf{11} (2019) 329--354; Dekel, Ely and Yilankaya,
\emph{The Review of Economic Studies} \textbf{74} (2007) 685--704; Heifetz,
Shannon and Spiegel, \emph{Journal of Economic Theory} \textbf{133} (2007)
31--57. Ours is the mirror image: their distortion is endogenous and sits in
behaviour, ours is an exogenous policy variable inside the selection
functional.

\textbf{Multilevel selection}: Wilson, \emph{Proceedings of the National
Academy of Sciences} \textbf{72} (1975) 143--146; Traulsen and Nowak,
\emph{Proceedings of the National Academy of Sciences} \textbf{103} (2006)
10952--10955. Ours is degenerate but sharper: one selecting level, not the
level that plays, selecting on a policy variable rather than on an aggregate of
the lower level.

\textbf{Multitask incentive distortion}: Holmstr\"{o}m, \emph{The Bell Journal
of Economics} \textbf{10} (1979) 74--91; Holmstr\"{o}m and Milgrom,
\emph{Journal of Law, Economics, and Organization} \textbf{7} (1991) 24--52;
Baker, \emph{Journal of Political Economy} \textbf{100} (1992) 598--614.
Liability is exactly such a distorted measure, taxing retaliation along with
initiation, but here the reallocation is evolutionary rather than deliberate,
and the loss is discontinuous in the power of the instrument rather than
smooth.

\textbf{Game-environment feedback}: Weitz et al., \emph{Proceedings of the
National Academy of Sciences} \textbf{113} (2016) E7518--E7525; Tilman, Plotkin
and Ak\c{c}ay, \emph{Nature Communications} \textbf{11} (2020) 915. Their
environmental state can recover and deployed capability cannot, which is what
makes our hysteresis width exactly computable. Nearby, Fern\'{a}ndez Domingos
et al., \emph{Scientific Reports} \textbf{12} (2022) 8492, and Terrucha et al.,
\emph{Scientific Reports} \textbf{14} (2024) 10460, treat delegation as a
commitment device chosen inside the game; our principal sits outside it.

\section*{4. Has any part been published before?}

No. No part of this paper has been published or presented before, at a
conference or elsewhere, and it is under review nowhere else. There is no
earlier version by the authors as a workshop paper, extended abstract or
preprint.

The one component we did not originate is the interaction layer, which is taken
unchanged from the framed AI-race experiment of Fern\'{a}ndez Domingos and Han
(preprint, arXiv:2607.26034, July 2026), by different authors: its stage game,
uncertain horizon, four reduced designs and parameter values are reused as
that preprint reports them. The reuse is deliberate, so that every result here
is attributable to the deployment layer alone; because every threshold is a
closed form in the payoff matrix, the results carry over unchanged in form to
any revised version of those constants. Computations use EGTtools
(Fern\'{a}ndez Domingos, Santos and Lenaerts, \emph{iScience} \textbf{26}
(2023) 106419). All code, tables and figures are public under the MIT licence
at \url{https://github.com/trungkiet2005/deployment-layer-selection}, and every
number quoted in the manuscript is regenerated by the released scripts and
pinned by its test suite.

\end{document}
